\documentclass[10pt,a4paper]{amsart}
\usepackage{amsmath,amssymb,amsthm,graphicx}
% For commutative diagrams we use Xy-pic
\usepackage[all]{xy}

% We use 2cell for 2-commutative diagrams.
\xyoption{2cell}
% Theorem environments.
%
\theoremstyle{plain}
\newtheorem{theorem}[subsection]{Theorem}
\newtheorem{proposition}[subsection]{Proposition}
\newtheorem{lemma}[subsection]{Lemma}

\theoremstyle{definition}
\newtheorem{definition}[subsection]{Definition}
\newtheorem{example}[subsection]{Example}
\newtheorem{exercise}[subsection]{Exercise}
\newtheorem{situation}[subsection]{Situation}

\theoremstyle{remark}
\newtheorem{remark}[subsection]{Remark}
\newtheorem{remarks}[subsection]{Remarks}

\numberwithin{equation}{subsection}

% Macros
%
\def\lim{\mathop{\mathrm{lim}}\nolimits}
\def\colim{\mathop{\mathrm{colim}}\nolimits}
\def\Spec{\mathop{\mathrm{Spec}}}
\def\Hom{\mathop{\mathrm{Hom}}\nolimits}
\def\Ext{\mathop{\mathrm{Ext}}\nolimits}
\def\SheafHom{\mathop{\mathcal{H}\!\mathit{om}}\nolimits}
\def\SheafExt{\mathop{\mathcal{E}\!\mathit{xt}}\nolimits}
\def\Sch{\mathit{Sch}}
\def\Mor{\mathop{\mathrm{Mor}}\nolimits}
\def\Ob{\mathop{\mathrm{Ob}}\nolimits}
\def\Sh{\mathop{\mathit{Sh}}\nolimits}
\def\NL{\mathop{N\!L}\nolimits}
\def\CH{\mathop{\mathrm{CH}}\nolimits}
\def\proetale{{pro\text{-}\acute{e}tale}}
\def\etale{{\acute{e}tale}}
\def\QCoh{\mathit{QCoh}}
\def\Ker{\mathop{\mathrm{Ker}}}
\def\Im{\mathop{\mathrm{Im}}}
\def\Coker{\mathop{\mathrm{Coker}}}
\def\Coim{\mathop{\mathrm{Coim}}}

% Boxtimes
%
\DeclareMathSymbol{\boxtimes}{\mathbin}{AMSa}{"02}

%
% Macros for moduli stacks/spaces
%
\def\QCohstack{\mathcal{QC}\!\mathit{oh}}
\def\Cohstack{\mathcal{C}\!\mathit{oh}}
\def\Spacesstack{\mathcal{S}\!\mathit{paces}}
\def\Quotfunctor{\mathrm{Quot}}
\def\Hilbfunctor{\mathrm{Hilb}}
\def\Curvesstack{\mathcal{C}\!\mathit{urves}}
\def\Polarizedstack{\mathcal{P}\!\mathit{olarized}}
\def\Complexesstack{\mathcal{C}\!\mathit{omplexes}}
% \Pic is the operator that assigns to X its picard group, usage \Pic(X)
% \Picardstack_{X/B} denotes the Picard stack of X over B
% \Picardfunctor_{X/B} denotes the Picard functor of X over B
\def\Pic{\mathop{\mathrm{Pic}}\nolimits}
\def\Picardstack{\mathcal{P}\!\mathit{ic}}
\def\Picardfunctor{\mathrm{Pic}}
\def\Deformationcategory{\mathcal{D}\!\mathit{ef}}
\usepackage{cleveref}
\crefname{originalreference}{original reference}{original references}
\Crefname{originalreference}{Original reference}{Original references}
\makeatletter
\newlabel{lemma-proper-base-change-f-star}{{0A3U}{source}{}{sourcepointer-1}{}}
\newlabel{lemma-proper-base-change-f-star@cref}{{[originalreference][0][]0A3U}{[1][0][]source}{}{}{}}
\makeatother
\crefname{originalreference}{original reference}{original references}
\Crefname{originalreference}{Original reference}{Original references}
\makeatletter
\newlabel{lemma-proper-base-change-in-terms-of-injectives}{{0A4B}{source}{}{sourcepointer-2}{}}
\newlabel{lemma-proper-base-change-in-terms-of-injectives@cref}{{[originalreference][0][]0A4B}{[1][0][]source}{}{}{}}
\makeatother
\crefname{originallemma}{original lemma}{original lemmas}
\Crefname{originallemma}{Original lemma}{Original lemmas}
\makeatletter
\newlabel{sites-cohomology-lemma-pushforward-injective-flat}{{0730}{source}{}{sourcepointer-3}{}}
\newlabel{sites-cohomology-lemma-pushforward-injective-flat@cref}{{[originallemma][0][]0730}{[1][0][]source}{}{}{}}
\makeatother
\crefname{originallemma}{original lemma}{original lemmas}
\Crefname{originallemma}{Original lemma}{Original lemmas}
\makeatletter
\newlabel{sites-cohomology-lemma-relative-Leray}{{0734}{source}{}{sourcepointer-4}{}}
\newlabel{sites-cohomology-lemma-relative-Leray@cref}{{[originallemma][0][]0734}{[1][0][]source}{}{}{}}
\makeatother
\title{Base change for a composite of proper morphisms follows from base change for the two factors}
\author{The Stacks Project Authors}
\date{stacks-project:etale-cohomology.tex:local-snapshot-20261011; GNU FDL 1.2 or later}
\begin{document}
\maketitle
\noindent\textit{Editorial note.} All delivered mathematical text is the author's own, copied from the pinned Stacks Project chapter file etale-cohomology.tex; no statement or proof step was rewritten, reordered or removed.\par\medskip
\noindent\textit{Editorial note.} The counted proof is the author's complete argument: the proof is written out in full: the reduction to an injective sheaf with injectivity of f\_*I, the base change square and the pushforward formula, the two vanishing inputs, and the degeneration of the relative Leray spectral sequence.\par\medskip
\noindent\textit{Editorial note.} The label lemma-composition names the counted statement, and the proof environment immediately following it in the source is delivered verbatim as its proof.\par\medskip
\noindent\textit{Editorial note.} The counted statement is carried by the author's own lemma environment in the chapter source.\par\medskip
\noindent\textit{Editorial note.} The statement and proof are copied from the pinned chapter file; the four cited inputs are kept as external pointers with their canonical Stacks tags.\par\medskip
\noindent\textit{Editorial note.} The preamble is the chapter's own native preamble excerpt already used for this source: the Xy-pic setup (native preamble lines 14-18) and the native theorem declarations with the mathematical macros (native preamble lines 154-217).\par\medskip
\noindent\textit{Allowed standard prerequisites (editorial).}\begin{itemize}
\item Base change reformulated through injective sheaves. Same chapter, the lemma identifying the base change property of a proper morphism with the vanishing of the higher direct images of injective torsion sheaves after base change. Supplies the form of the hypothesis and the form of the conclusion.
\item Pushforward of an injective sheaf along a flat/proper morphism. Cohomology on Sites, the lemma making f\_*I injective for I injective and f flat in the relevant sense. Supplies injectivity of f\_*I, which lets the hypothesis on g be applied to it.
\item The relative Leray spectral sequence and the pushforward base change formula. Cohomology on Sites, the relative Leray spectral sequence for a composite, and Same chapter, the pushforward base change lemma b\textasciicircum{}\{-1\}f\_*I = f'\_*a\textasciicircum{}\{-1\}I. Supplies the spectral sequence whose degeneration gives the conclusion and the identification of its middle term.
\end{itemize}\medskip
\ifdefined\hypertarget\hypertarget{sourcepointer-1}{}\fi
\noindent\textit{Source reference.} Original reference 0A3U: sources/stacks-etale-cohomology/pinned-original-source.tex, source lines 18282--18282. Prerequisite pointer invoked at proof line(s) 18504; it names a standard fact of the same chapter that the counted argument uses, and every step of the counted proof built on it is retained verbatim in the delivered body.\par\medskip
\ifdefined\hypertarget\hypertarget{sourcepointer-2}{}\fi
\noindent\textit{Source reference.} Original reference 0A4B: sources/stacks-etale-cohomology/pinned-original-source.tex, source lines 18360--18360. Prerequisite pointer invoked at proof line(s) 18488; it names a standard fact of the same chapter that the counted argument uses, and every step of the counted proof built on it is retained verbatim in the delivered body.\par\medskip
\ifdefined\hypertarget\hypertarget{sourcepointer-3}{}\fi
\noindent\textit{Source reference.} Original lemma 0730: sources/stacks-chow/tags.txt, source lines 9044--9044. Prerequisite pointer invoked at proof line(s) 18495; it names a standard fact of a sibling chapter that the counted argument uses, and every step of the counted proof built on it is retained verbatim in the delivered body.\par\medskip
\ifdefined\hypertarget\hypertarget{sourcepointer-4}{}\fi
\noindent\textit{Source reference.} Original lemma 0734: sources/stacks-chow/tags.txt, source lines 9048--9048. Prerequisite pointer invoked at proof line(s) 18509; it names a standard fact of a sibling chapter that the counted argument uses, and every step of the counted proof built on it is retained verbatim in the delivered body.\par\medskip
\begin{lemma}
\label{lemma-composition}
Let $f : X \to Y$ and $g : Y \to Z$ be proper morphisms of schemes. Assume
\begin{enumerate}
\item cohomology commutes with base change for $f$, and
\item cohomology commutes with base change for $g$.
\end{enumerate}
Then cohomology commutes with base change for $g \circ f$.
\end{lemma}

\begin{proof}
We will use the equivalence of
Lemma \ref{lemma-proper-base-change-in-terms-of-injectives}
without further mention. Let $\ell$ be a prime number.
Let $\mathcal{I}$ be an injective sheaf of
$\mathbf{Z}/\ell\mathbf{Z}$-modules on $X_\etale$.
Then $f_*\mathcal{I}$ is an injective sheaf of
$\mathbf{Z}/\ell\mathbf{Z}$-modules on $Y_\etale$
(Cohomology on Sites, Lemma
\ref{sites-cohomology-lemma-pushforward-injective-flat}).
The result follows formally from this, but we will also
spell it out.

\medskip\noindent
Let $Z' \to Z$ be a morphism of schemes and set
$Y' = Z' \times_Z Y$ and $X' = Z' \times_Z X = Y' \times_ Y X$.
Denote $a : X' \to X$, $b : Y' \to Y$, and $c : Z' \to Z$ the
projections. Similarly for $f' : X' \to Y'$ and $g' : Y' \to Z'$.
By Lemma \ref{lemma-proper-base-change-f-star} we have
$b^{-1}f_*\mathcal{I} = f'_*a^{-1}\mathcal{I}$.
On the other hand, we know that $R^qf'_*a^{-1}\mathcal{I}$ and
$R^q(g')_*b^{-1}f_*\mathcal{I}$ are zero for $q > 0$.
Using the spectral sequence
(Cohomology on Sites, Lemma \ref{sites-cohomology-lemma-relative-Leray})
$$
R^pg'_* R^qf'_* a^{-1}\mathcal{I} \Rightarrow
R^{p + q}(g' \circ f')_* a^{-1}\mathcal{I}
$$
we conclude that $R^p(g' \circ f')_*a^{-1}\mathcal{I} = 0$ for
$p > 0$ as desired.
\end{proof}

\end{document}
